\section{Motivation}
\label{intro:sec:motivation}

Artificial Intelligence(AI) is the buzzword of this decade.
The release of ChatGPT in late 2022 triggered a wave of public interest in machine
learning that had been building for years in research laboratories.
What made this transition possible was the solution of an extraordinary engineering
challenge: the construction of a pipeline capable of digesting nearly the entirety
of human digital knowledge and using it to train deep learning models of
unprecedented scale.
This achievement draws on decades of cumulative work spanning hardware design,
software engineering, and fundamental research: from the first attempts to emulate
the human brain with artificial neural networks in the 1950s, through the
development of backpropagation in the 1980s, to the rise of GPU computing in the
2000s and the explosion of large-scale datasets in the 2010s.
Yet for all this engineering progress, our theoretical understanding of \emph{why}
these models work remains strikingly limited.

The reasons to close this gap are manifold.
The most immediate are practical: a deeper understanding can guide the design of
better algorithms, improve the efficiency of existing models, and reduce the
enormous computational and energetic costs of training.
The social motivations are equally compelling: a rigorous theory underpins the
growing fields of interpretability, fairness, and robustness, which are essential
if AI systems are to be deployed safely in high-stakes settings.
But this thesis is submitted to the \emph{École Doctorale de Physique} (EDPY), and
physicists are not driven by practical concerns alone.
We are also moved by the desire to understand the fundamental principles governing
the behavior of complex systems.
There is something deeply human in wanting to understand the mechanisms of a
technology that increasingly surrounds us and shapes every aspect of our lives;
satisfying that curiosity is a legitimate scientific goal in its own right, and has
historically led to unexpected discoveries with far-reaching consequences.

What makes deep learning theoretically challenging is precisely what makes it
empirically powerful.
Many machine learning models are theoretically grounded by construction: support
vector machines rest on the geometry of convex optimization; decision trees can be
analyzed with combinatorial and probabilistic tools.
Deep learning broke free of these clean frameworks.
The central idea, using a very large number of parameters to fit data, violates
the assumptions underlying classical statistics: the bias--variance trade-off,
convexity, and the availability of closed-form solutions.
Yet deep networks learn effective representations, generalize remarkably well to
unseen examples, and are trained with algorithms as simple as stochastic gradient
descent navigating a highly non-convex loss landscape.
This has challenged many of the intuitions we had about learning and generalization,
and continues to generate vigorous debate in the research community.

A physicist's instinct, confronted with a phenomenon that outstrips existing theory,
is to build a model, a simplified but principled caricature that retains the
essential features and can be understood analytically.
Think of where chemistry stood at the end of the nineteenth century: practitioners
could synthesize remarkably complex molecules and predict reactions with
considerable reliability, yet nobody understood the deep mechanism, why atoms bond
the way they do.
The Thomson model was the first crude, almost certainly incomplete picture of the
underlying atomic structure.
It was wrong in important ways, but concrete enough to be tested, refuted, and
improved.
This is, roughly speaking, where we stand with deep learning today, and it is the
role we aspire to play.
We want to be Thomson, so that someone can be Bohr after us, and eventually
someone arrives at quantum mechanics.

One advantage we have is that our experiments run on computers.
We do not need a particle collider to test a theory of deep learning: we need a
laptop (or a cluster for large-scale experiments) and a few days of compute time.
This is an enormous freedom, and we should exploit it.
A second advantage is that the fundamental process at the core is mathematically
well-defined: training a neural network is an optimization procedure whose rules are
written down precisely, even if their consequences are not.
We do not need to figure out \emph{what} the process is; only \emph{why} it works
and \emph{how} its dynamics unfold.
The goal is a theory that is both analytically tractable and grounded in
experimental evidence.

Among the many aspects of deep learning that resist theoretical explanation, we
focus on the \emph{dynamics of training}: the process by which randomly initialized
parameters converge to a configuration that generalizes well to unseen data.
The loss landscape of deep learning is highly non-convex, and it is far from obvious
why stochastic gradient descent should find good solutions in it.
Yet it does, and the dynamics induced by SGD introduce an inductive bias that
correlates with generalization.
Understanding this bias, its origins, and its consequences is the central thread
running through this thesis.

Classical learning theory---PAC learning, VC dimension, Rademacher complexity and all this jargon of tools---analyzes generalization in the worst case over all possible data
distributions.
This regime is useful as a foundation but tends to produce pessimistic bounds that
badly overestimate the true error of well-trained deep networks.
Statistical physics offers a complementary perspective, analyzing the \emph{typical}
behavior of learning algorithms averaged over a natural distribution of instances.
This leads to sharper, more realistic predictions.
The connection between statistical physics and machine learning is not new: the
tools of disordered systems and mean-field theory were applied to neural network
capacity and generalization as far back as the 1980s.
In some sense, machine learning is a canonical problem of statistical physics: large
numbers of parameters, large datasets, and emergent collective behavior.
We inherit and build on this tradition.
